Un astronome amateur observe la Lune, située à \qty{3.84e8}{m} de la Terre, avec un télescope dont l'objectif a un diamètre de \qty{\v{Dt}}{cm}. On considère une lumière de longueur d'onde \qty{\v{lam}}{nm} et on néglige les effets de l'atmosphère.

\begin{enumerate}[label=\alph*)]
    \item Déterminez l'angle limite de résolution du télescope.
    \item Quelle est la plus petite distance entre deux détails de la surface lunaire que ce télescope peut distinguer ?
    \item Quel diamètre d'objectif faudrait-il pour distinguer deux détails séparés de \qty{\v{s}}{m} ?
\end{enumerate}
